\section{Controller-laget}
\label{sec:controllers}

API-et har to kontrollere: \texttt{RoomsController} og
\texttt{AuthController}. Begge bruker \texttt{[ApiController]}, som gir
automatisk modellvalidering og standardiserte feilsvar. Kontrollerne er
tynne: de binder forespørsler, delegerer til provider-laget og logger
vellykkede operasjoner, men inneholder ingen forretningslogikk og ingen
\texttt{try/catch}-blokker.

\texttt{RoomsController} er beskyttet med \texttt{[Authorize]} på
klassenivå og eksponerer syv endepunkter under \texttt{/api/rooms} som
dekker kravene F1.1--F1.5 og S1.2 \cite{hdo_kravspec}. Riktig
videomotor velges via et prefiks kodet inn i ressurs-IDen: ved oppretting
sendes ønsket provider i request-bodyen og bakes inn i den returnerte
IDen på formen \texttt{provider:rawId}. På øvrige endepunkter parses
prefikset ut av \texttt{roomId} og brukes til å slå opp riktig
\texttt{IVideoProvider} via \texttt{IVideoProviderFactory}.
Feilhåndtering er sentralisert i \texttt{GlobalExceptionHandler}, som
mapper typede unntak fra provider-laget til HTTP-statuskoder:
\texttt{NotFoundException} gir 404, \texttt{NotSupportedException} gir
501 og \texttt{ExternalServiceException} gir 502. Alle feilsvar følger
RFC~7807 ProblemDetails-formatet.

\texttt{AuthController} har ett endepunkt, \texttt{POST /auth/token},
som henter et Bearer-token via OAuth~2.0 client credentials-flyten.
Endepunktet er annotert med \texttt{[AllowAnonymous]} siden det må nås
før et token er tilgjengelig, og med \texttt{[EnableRateLimiting("auth")]}
for å begrense antall forsøk per IP. Selve kommunikasjonen mot Keycloak
er delegert til \texttt{KeycloakService}, som POSTer
\texttt{client\_id} og \texttt{client\_secret} til Keycloaks
\texttt{/protocol/openid-connect/token}-endepunkt og returnerer
access-tokenen. Tjenesten er stateless og uten caching — ansvaret for å
gjenbruke tokenen ligger hos kalleren.